\section{Nonadic connection, calendar extension and transport memory}
\label{cap:calendario}

TPK chronology combines an observable phase, a dodecaphasic register
and a memory of the transitions performed. Their relation is determined by
the composition of displacements. Crossing the phase boundary
increments a quotient; preserving that quotient distinguishes
successive returns. This structure provides the arithmetic basis of
nonadic holonomy and its bilateral representation.

It is useful to fix from the outset the three objects involved. The phase
belongs to a cycle of nine positions. The finite calendar retains twelve
turns of that cycle. The enriched history additionally preserves the integer
number of turns and the registers of position, orientation, sheet, carry and
boundary. Each reduction preserves a specific observation of the same
evolution; the available information is specified through its projections.

\subsection{Incidence of arithmetic modes}
\label{sec:cal-incidencia}

\subsubsection{Selection, persistence and cyclic succession}

Let $\Sigma$ and $\Pi$ be the modes corresponding to the additive
and multiplicative APP evaluations. In each block of three
events, the phase selector traverses the sequence $(+1,-1,0)$. Its action on mode $m$ is
\[
 +1:m\longmapsto\Sigma,\qquad
 -1:m\longmapsto\Pi,\qquad
 0:m\longmapsto m.
\]
The neutral event preserves the preceding mode. Starting from the cyclic closure
of the block, the three recorded modes are therefore
\[
 (\Sigma,\Pi,\Pi).
\]
The third position preserves the multiplicative-mode label
selected in the preceding event. In the emitter, that neutral event
increments only the counter $Z$. This distinction allows
the incidence to be constructed directly from the calendar transitions.

\begin{proposicion}[Modal incidence of the tritic block]
\label{prop:cal-incidencia}
Counting the consecutive pairs of the block, including its closure, and ordering
the modes as $(\Sigma,\Pi)$ gives
\begin{equation}
 F_{\rm av}=\begin{pmatrix}0&1\\1&1\end{pmatrix}.
 \label{eq:cal-incidencia}
\end{equation}
In a sequence of $3q$ events consisting of $q$ repetitions of the block,
the transition count is $qF_{\rm av}$.
\end{proposicion}
\begin{proof}
The ordered pairs are $\Sigma\to\Pi$, $\Pi\to\Pi$ and
$\Pi\to\Sigma$. Each appears once; $\Sigma\to\Sigma$
appears zero times. These are, respectively, the upper-right,
lower-right, lower-left and upper-left entries of the
matrix. Each repetition adds exactly the same three pairs, including
the link from the last mode of one block to the first of the next.
The sum of the counts is $qF_{\rm av}$.
\end{proof}

The lengths $9$, $27$ and $108$ thus produce $3F_{\rm av}$,
$9F_{\rm av}$ and $36F_{\rm av}$. These matrices count events of an
already fixed periodic itinerary. The powers of $F_{\rm av}$ count
concatenations admitted by its incidence. These are two distinct
mathematical operations on the same matrix: summation of occurrences and
composition of transitions.

\subsubsection{Concatenation of incidences and induced recurrence}

Let us define integers $a_0=0$, $a_1=1$ and
$a_{n+1}=a_n+a_{n-1}$ for $n\geq1$. Their appearance follows from
multiplication of the matrix already constructed.

\begin{proposicion}[Powers of modal incidence]
\label{prop:cal-potencias}
For every $n\geq1$,
\[
 F_{\rm av}^{\,n}=
 \begin{pmatrix}a_{n-1}&a_n\\a_n&a_{n+1}\end{pmatrix},
 \qquad
 a_{n-1}a_{n+1}-a_n^2=(-1)^n.
\]
\end{proposicion}
\begin{proof}
For $n=1$, the formula reproduces \eqref{eq:cal-incidencia} because
$a_2=1$. If the formula holds for $n$, multiplication by $F_{\rm av}$
gives
\[
 \begin{pmatrix}
 a_n&a_{n-1}+a_n\\a_{n+1}&a_n+a_{n+1}
 \end{pmatrix}
 =\begin{pmatrix}a_n&a_{n+1}\\a_{n+1}&a_{n+2}\end{pmatrix}.
\]
Induction proves the first identity. Taking determinants and using
$\det F_{\rm av}=-1$ gives the second.
\end{proof}

In particular, $F_{\rm av}^{2}=F_{\rm av}+I$. This identity allows
a hyperbolic generator to be normalized through
\begin{equation}
 J_{-1}=\frac{2F_{\rm av}^{2}-3I}{\sqrt5}
       =\frac{2F_{\rm av}-I}{\sqrt5},
 \qquad J_{-1}^{2}=I.
 \label{eq:cal-generador-hiperbolico}
\end{equation}
Indeed, $(2F_{\rm av}-I)^2=4(F_{\rm av}+I)-4F_{\rm av}+I=5I$.
The coefficient $5$ comes here from the matrix identity. The quadratic
representation and Euler extension use this generator after its
modal construction.

The recurrence $(a_n)$ subsequently admits identification with the
Fibonacci sequence. The proof order in this subsection is explicit:
tritic selector, sequence of modes, edge count, incidence and
recurrence. Identification of the sequence preserves that order.
Coinductive regional generation also has its own emitters, filters,
lifts and survival conditions; the modal recurrence specifies
one of its structural realizations.

\subsection{Arithmetic extension of the calendar}

\subsubsection{Decomposition into phase and dodecaphasic position}

We write $C_q=\ZZ/q\ZZ$ for the additive group of residues modulo $q$.
The position of a calendar of $108=12\cdot9$ events has a
unique decomposition
\[
 t=9b+r,\qquad 0\leq b<12,\quad0\leq r<9,
\]
for the representative $0\leq t<108$. The coordinate $r$ is the nonadic
phase; $b$ identifies the position of the turn in the dodecaphasic
register. The scale $108$ brings those two registers together according to the law of
composition with carry.

\begin{teorema}[Calendar extension with carry]
\label{thm:cal-extension}
The maps
\[
 \iota:C_{12}\longrightarrow C_{108},\quad[b]\longmapsto[9b],
 \qquad p:C_{108}\longrightarrow C_9,\quad[t]\longmapsto[t]_9
\]
determine an exact sequence
\begin{equation}
 0\longrightarrow C_{12}\xrightarrow{\ \iota\ }C_{108}
 \xrightarrow{\ p\ }C_9\longrightarrow0.
 \label{eq:cal-extension}
\end{equation}
In the coordinates $(b,r)$, addition is given by
\begin{equation}
 (b,r)\oplus(b',r')=
 \left(\left[b+b'+\left\lfloor\frac{r+r'}9\right\rfloor\right]_{12},
 [r+r']_9\right).
 \label{eq:cal-ley-grupo}
\end{equation}
This extension has no homomorphic section.
\end{teorema}
\begin{proof}
The equality $[9b]=0$ in $C_{108}$ is equivalent to $12\mid b$, so
$\iota$ is injective. The map $p$ is surjective and its kernel
consists exactly of the classes that are multiples of nine, the image of
$\iota$. This demonstrates exactness.

To add two representatives, one writes
\[
 9b+r+9b'+r'
 =9\left(b+b'+\left\lfloor\frac{r+r'}9\right\rfloor\right)
   +[r+r']_9.
\]
Reducing the first coefficient modulo twelve gives
\eqref{eq:cal-ley-grupo}. Now suppose that a homomorphic
section $s:C_9\to C_{108}$ exists. The representative $a$ of $s([1])$
satisfies $a\equiv1\pmod9$ and $9a\equiv0\pmod{108}$. The second
condition implies $12\mid a$, whence $a\equiv0\pmod3$; the first
implies $a\equiv1\pmod3$. The contradiction demonstrates the last
assertion.
\end{proof}

\subsubsection{Cocycle, inverse and restitution}

The term
\[
 \kappa_9(r,s)=\left\lfloor\frac{r+s}9\right\rfloor,
 \qquad 0\leq r,s<9,
\]
records the threshold crossing. Its composition identity is
\begin{equation}
 \kappa_9(r,s)+\kappa_9([r+s]_9,u)
 =\kappa_9(s,u)+\kappa_9(r,[s+u]_9).
 \label{eq:cal-cociclo}
\end{equation}
Both sides are $\lfloor(r+s+u)/9\rfloor$: it suffices to substitute
$r+s=9\kappa_9(r,s)+[r+s]_9$, or the analogous decomposition of $s+u$.
The identity ensures that the accumulated quotient is independent of the
parenthesization of three summands.

The inverse of $(b,r)$ is determined by
\[
 -(b,r)=
 \begin{cases}
 ([-b]_{12},0),&r=0,\\
 ([-b-1]_{12},9-r),&1\leq r<9.
 \end{cases}
\]
The additional term $-1$ represents the exact boundary correction when
inverting an interior residue. For example, the representative $t=17$ has
coordinates $(1,8)$. Its inverse is $(-2,1)$ in coordinates with
first component modulo twelve: $9(-2)+1=-17$. The inverted residue and
the corrected quotient jointly restore the opposite integer.

\subsection{Integer lift and nonadic holonomy}

\subsubsection{Finite phase and turn counter}

Euclidean division extends to every $k\in\ZZ$:
\[
 k=9m+r,\qquad m=\left\lfloor\frac k9\right\rfloor,
 \quad0\leq r<9.
\]
The pair $(m,r)$ determines $k$ exactly. Its reduction
$(m,r)\mapsto([m]_{12},r)$ recovers the finite calendar. In the
integer lift, a turn of nine steps acts by
\begin{equation}
 (r,m)\longmapsto(r,m+1).
 \label{eq:cal-vuelta-entera}
\end{equation}
Twelve turns restore both finite coordinates and shift the
integer counter by twelve units. Enriched memory includes that
counter together with the ordered register of transitions.

\subsubsection{Composition of admissible prolongations}

At level $k$, we represent the enriched regional state by
\[
 \widehat x_k=(B_k,C_k,p_k,c_k,\Sigma_k,W_k,g_k,m_k).
\]
$B_k$ brings together the regional blocks; $C_k$ is the compatible cylinder;
$p_k$ is the prefix; $c_k$ records the carry; $\Sigma_k$ preserves the
boundary and $W_k$ the incidence. The last coordinates are the phase and
the turn counter. The prolongation operators act on
this complete state, with the dependencies of each component.

\begin{definicion}[Nonadic holonomy of the prolongation system]
Let $\mathcal R_k\subseteq\widehat X_k\times\widehat X_{k+1}$ be
the admissible relations produced by the TPK. The composition of one
turn is
\[
 \Gamma_{9,k}=\mathcal R_{k+8}\circ\cdots\circ\mathcal R_k.
\]
For $s\geq1$, the composition of $s$ turns is
\[
 \Gamma_{9s,k}=
 \Gamma_{9,k+9(s-1)}\circ\cdots\circ\Gamma_{9,k}.
\]
\end{definicion}

Relational composition retains the prolongations that simultaneously
satisfy the constraints of the state. An admissible trajectory
singles out a sequence of such transitions. The notation distinguishes
the admissibility relation, the chosen trajectory and the finite
observation of its phase.

\begin{proposicion}[Phase return and memory advance]
\label{prop:cal-memoria}
On a trajectory whose turn rule is
$g_{k+9}=g_k$, $m_{k+9}=m_k+1$, every admissible composition of $s$
turns satisfies
\[
 g_{k+9s}=g_k,\qquad m_{k+9s}=m_k+s,
 \qquad [m_{k+9s}]_{12}=[m_k+s]_{12}.
\]
\end{proposicion}
\begin{proof}
The case $s=1$ is the turn rule. Composing an additional turn
with the trajectory already constructed preserves the phase and the counter
receives another unit. Induction demonstrates the first two equalities;
projection modulo twelve gives the third.
\end{proof}

The register allows each return to be recognized as a transformation of
history. The observable phase expresses a periodicity, while
the counter and boundary coordinates distinguish its successive
realizations. The regional prolongation proofs additionally specify
which cylinders survive and which prefixes are determined.

\subsection{Bilateral representation of phase transport}

\subsubsection{Unit shift and group of turns}

The lift of the index admits a representation in
$\mathcal H_{\rm ind}=\ell^2(\ZZ)$. Its orthonormal basis is written
$e_k=e_{r,m}$ when $k=9m+r$. We define $Te_k=e_{k+1}$, that is,
\[
 Te_{r,m}=\begin{cases}
 e_{r+1,m},&r<8,\\e_{0,m+1},&r=8.
 \end{cases}
\]
The transformation bijectively permutes an orthonormal basis and
extends by continuity to a unitary operator. The inverse takes
$e_{r,m}$ to $e_{r-1,m}$ when $r>0$ and takes $e_{0,m}$ to
$e_{8,m-1}$.

\begin{teorema}[Faithful representation of the turn counter]
\label{thm:cal-representacion}
Let $S=T^9$. Then
\[
 S^se_{r,m}=e_{r,m+s},\qquad s\in\ZZ.
\]
The map $s\mapsto S^s$ is a faithful unitary representation of
$(\ZZ,+)$.
\end{teorema}
\begin{proof}
Nine steps preserve the residue and increase the quotient by one unit.
The same identity applied to the inverse provides the formula for
negative exponents. The powers satisfy
$S^{s+t}=S^sS^t$. If $S^s=I$, the equality
$e_{r,m+s}=e_{r,m}$ forces $s=0$ by independence of the basis
vectors. This proves faithfulness.
\end{proof}

The index space is a realization of the counter. The TPK state
space also incorporates the arithmetic and incidence data.
Unitary inversion of the index applies to its bilateral covering;
inversion of a complete transition uses the archive of that
transition and its admissibility conditions.

\subsubsection{Memory operator and boundary correction}

Let us define, on finitely supported vectors,
\[
 D_Me_k=\left\lfloor\frac k9\right\rfloor e_k,
 \qquad \mathcal Ie_k=e_{-k},
 \qquad P_0e_k=\begin{cases}e_k,&9\mid k,\\0,&9\nmid k.\end{cases}
\]
This domain is dense and stable under the operations considered.

\begin{proposicion}[Memory and inversion identities]
On that domain the following hold
\begin{equation}
 [D_M,S]=S,\qquad
 \mathcal IT\mathcal I=T^{-1},\qquad
 \mathcal ID_M\mathcal I=-D_M-(I-P_0).
 \label{eq:cal-inversion}
\end{equation}
\end{proposicion}
\begin{proof}
If $k=9m+r$, then
$D_MS e_k=(m+1)e_{k+9}$ and $SD_Me_k=me_{k+9}$, which proves
the commutator. The composition $\mathcal IT\mathcal I$ takes
$e_k$ to $e_{k-1}$. Finally,
\[
 \left\lfloor\frac{-k}9\right\rfloor
 =\begin{cases}-m,&r=0,\\-m-1,&r>0.\end{cases}
\]
The diagonal action of the last member of \eqref{eq:cal-inversion}
coincides with that expression on each basis vector. Linearity
concludes the proof.
\end{proof}

The correction $I-P_0$ locates the boundary contribution in the interior
phases. When the path is inverted, the quotient and residue are
transformed jointly. That identity expresses arithmetically
the information that disappears when only the phase is observed.

\subsection{Refinement capacity and arithmetic vacancies}

\subsubsection{Exact comparison of ternary and decimal blocks}

A block of six trits has $3^6=729$ configurations in its
ambient space. A decimal block of three digits has $10^3=1000$
positions. Their cardinalities define an arithmetic encoding
capacity. For $n\geq0$, let
\begin{equation}
 \mathcal C_n=\max\{a\in\NN:1000^a\leq729^{n+1}\}.
 \label{eq:cal-capacidad-entera}
\end{equation}
The definition compares integers and can be evaluated exactly. In a
subsequent logarithmic representation,
\[
 \rho=\log_{1000}729=\log_{10}9,\qquad
 \delta=1-\rho,\qquad
 \mathcal C_n=\lfloor(n+1)\rho\rfloor.
\]
The ratio $\rho$ belongs to $(0,1)$ and is irrational. If
$\rho=p/q$ with positive integers, one would have $10^p=9^q$;
the prime factorization of the two sides prevents this.

\begin{definicion}[Register of capacity vacancies]
We define
\[
 Z_n=\lfloor(n+1)\delta\rfloor,\qquad
 z_n=\lfloor(n+1)\delta\rfloor-\lfloor n\delta\rfloor,
 \quad n\geq0.
\]
The value $z_n\in\{0,1\}$ records capacity retention at the
corresponding step. One has $z_0=0$ at the chosen origin.
\end{definicion}

\begin{teorema}[Balance of capacity and vacancies]
\label{thm:cal-capacidad}
For $n\geq0$ in the first identity and $n\geq1$ in the second,
\begin{equation}
 \mathcal C_n+Z_n=n,
 \qquad \mathcal C_n-\mathcal C_{n-1}=1-z_n.
 \label{eq:cal-capacidad}
\end{equation}
In particular, for $0\leq a<b$,
\[
 \mathcal C_b-\mathcal C_a+\sum_{n=a+1}^{b}z_n=b-a.
\]
\end{teorema}
\begin{proof}
The irrationality of $\delta$ implies
$\lceil(n+1)\delta\rceil=\lfloor(n+1)\delta\rfloor+1$.
Therefore,
\[
 \lfloor(n+1)(1-\delta)\rfloor
 =n+1-\lceil(n+1)\delta\rceil=n-Z_n.
\]
The difference between two consecutive equalities gives the second identity
for $n\geq1$. The telescoping sum demonstrates the formula on an interval.
\end{proof}

The register $Z_n$ can be obtained through $n-\mathcal C_n$,
using only the integer comparison
\eqref{eq:cal-capacidad-entera}. The logarithmic functions express
the same construction in analytic coordinates. The first capacity
retention occurs between $n=20$ and $n=21$:
\[
 \mathcal C_{19}=19,\quad\mathcal C_{20}=20,\quad
 \mathcal C_{21}=20,\quad\mathcal C_{22}=21.
\]
The history receives an additional block while decimal capacity
remains constant. The vacancy records precisely that offset.

\subsubsection{Distinguishing the operational indices}

The phase of the finite emitter, the regional prolongation phase and decimal
capacity belong to declared indices. The nonadic connection preserves
the phase after nine steps of its own index; vacancy preserves a
capacity during one step of the block index. The comparison between
the two is expressed by
\[
 [\mathcal C_n]_q=\bigl[[n]_q-[Z_n]_q\bigr]_q,
 \qquad q=9\ \text{or}\ 108.
\]
The phase offset depends on the accumulated vacancy register. Preserving it
allows one phase reading to be transformed into the other without improperly
identifying the two calendars.

The resulting structure links modal selection, carry composition,
holonomy and refinement. Memory fulfills precise functions at each
level: it restores quotients, distinguishes turns, records transitions and
controls encoding capacity. These operations constitute the
temporal and archival support on which regional readers
and realizations of unit conservation act.
